% =====================================================================
%  第 2 章  已有工作
% =====================================================================

\section{已有工作}
\label{sec:background}

本节简要交代与本课题直接相关的工作，以说明本课题的位置，不展开技术细节。

\subsection{自动猜想程序}

最早系统性地从事该项工作的是 Fajtlowicz 于二十世纪八十年代开发的 Graffiti。该程序建立一个图数据库，为每张图计算各种不变量，再以一组人工设计的启发式规则挑选可能成立的不等式。据其作者统计，Graffiti 累计提出过上千条猜想，其中一部分后来被证明或推翻。

此后出现了若干改进版本，较知名的有 DeLaViña 的 Graffiti.pc、Caporossi 与 Hansen 的 AutoGraphiX，以及近年 Davila 的 TxGraffiti。这些工作在两个方向上有所加强：以线性规划代替手工规则确定系数，以及主动搜索反例而非被动检验。

上述系统的共同做法是产出候选、交由人工挑选。它们能够生成大量不等式，但均未将“如何从上千条候选中可靠地筛出少数几条”作为核心问题处理。

\subsection{基于大模型的方法}

2024 年以来，出现了将大语言模型接入数学发现流程的工作，代表性成果为 FunSearch 与 AlphaEvolve，以及 Charton 的 PatternBoost。这些工作在组合构造问题上取得了超过此前记录的结果。

这些工作处理的是构造问题而非猜想问题：其目标是找到一个具体的对象，而不是提出一条对全部对象成立的普遍命题。构造问题的优劣可由目标函数直接衡量，因而不需要处理候选过多、真假难辨的筛选问题。

本课题关注的正是后者。表\ref{tab:position} 给出两类工作与本课题的对比。

\begin{table}[H]
\centering
\small
\caption{相关工作与本课题的定位}
\label{tab:position}
\begin{tabular}{@{}lcc@{}}
\toprule
 & 是否处理候选过多 & 是否使用大语言模型 \\
\midrule
Graffiti 系列 & 否（人工挑选） & 否 \\
FunSearch 系列 & 不适用（构造型任务） & 是 \\
本课题 & 是（统计筛选） & 是 \\
\bottomrule
\end{tabular}
\end{table}

需要指出，上述空缺并非理论上的深刻发现，而更接近于一个被忽略的工程问题：研究多集中于提升生成的数量与质量，较少讨论生成之后如何处理。本课题的价值不在于提出新理论，而在于将这一问题认真处理一次，并给出可运行的实现与初步实证。
